% !TEX program = xelatex
\documentclass[reprint,amsmath,amssymb,aps,pra,10pt]{revtex4-2}

\usepackage{fontspec}
\setmainfont{Arial}
\setsansfont{Arial}
\renewcommand{\familydefault}{\sfdefault}
\usepackage{graphicx}
\usepackage{bm}

\begin{document}

\title{Learning Equilibrium Polymer Distributions with Continuous Normalizing Flows and Diffusion Models}

\author{Author Name}
\affiliation{Affiliation}

\date{\today}

\begin{abstract}
Our long-term goal is to learn models that can support molecular trajectory inference from static structural samples, including protein conformational ensembles.
We first study controlled Brownian-dynamics synthetic systems where we can evaluate the equilibrium distribution and score field directly.
The codebase uses Rouse-chain simulations as the main data source and compares two likelihood- or score-based models: an equivariant continuous normalizing flow (CNF) and an equivariant variance-preserving diffusion model.
We train both models on centered coordinate snapshots and evaluate them by generated-configuration quality and agreement between learned log-density gradients and Brownian drift forces.
\end{abstract}

\maketitle

\section{Motivation}

Our long-term goal is to develop generative models that can help infer molecular conformational dynamics from static structural samples.
This problem is especially relevant for proteins, where experimental methods can provide structural ensembles but often do not provide the time ordering or transition information needed to reconstruct dynamics directly.
Atomistic proteins make this idea difficult to evaluate directly because experiments do not fully observe the effective dynamics and do not provide the correct score or force field as a clean target.

We therefore begin with controlled Brownian-dynamics synthetic systems.
In this setting, a known stochastic process generates equilibrium samples, and the simulator provides the corresponding score field for evaluation.
This lets us test whether continuous normalizing flows and diffusion models can learn both the equilibrium distribution of configurations and the local log-density gradients that govern overdamped dynamics.

For a configuration $x_t \in \mathbb{R}^{d \times N}$ with $N$ beads in $d$ spatial dimensions, we simulate the dynamics
\begin{equation}
    dX_t = F(X_t)\,dt + \sqrt{2D}\,dW_t,
\end{equation}
where $D$ denotes the diffusion coefficient, $W_t$ denotes standard Brownian motion, and $F(x)$ denotes the deterministic drift.
For the current base Rouse simulations, $D=0.25$.
For the Rouse-chain experiments, adjacent harmonic springs generate the base drift,
\begin{equation}
    F_i(x)
    =
    k_{\xi}
    \left(
        x_{i+1} - 2x_i + x_{i-1}
    \right),
\end{equation}
with endpoint modifications.
The implementation also supports nonideal terms, including Lennard-Jones interactions, excluded volume, confinement, hairpin contacts, and ring closure.

At equilibrium, the stationary-density score determines the drift:
\begin{equation}
    F(x) = D \nabla_x \log p_{\mathrm{eq}}(x).
\end{equation}
This relation gives a practical test for learned generative models.
If a model learns $p_{\mathrm{eq}}(x)$, then its learned score should align with the Brownian drift divided by $D$.
For the current base Rouse setting, the score target is therefore $F(x)/0.25 = 4F(x)$.
The current experiments therefore learn equilibrium coordinate distributions from static snapshots, then compare generated samples and learned score fields against known Rouse-chain quantities.

\section{Hamiltonian and Boltzmann Background}

For atomistic molecular systems, we usually write the longer-term trajectory-inference framing in Hamiltonian form.
Let $\vec{q}$ denote particle positions and $\vec{p}$ denote momenta.
Hamilton's equations take the form
\begin{align}
    \frac{dq_{i,k}}{dt}
    &=
    \frac{\partial H}{\partial p_{i,k}},
    \\
    \frac{dp_{i,k}}{dt}
    &=
    -
    \frac{\partial H}{\partial q_{i,k}},
\end{align}
where $i$ indexes particles and $k$ indexes spatial coordinates.
For standard molecular dynamics, momenta satisfy
\begin{equation}
    p_{i,k} = m_i \frac{dq_{i,k}}{dt},
\end{equation}
and we decompose the Hamiltonian into kinetic and potential energy:
\begin{equation}
    H(\vec{q},\vec{p},t)
    =
    KE(\vec{p}) + PE(\vec{q})
    =
    \sum_i \sum_k \frac{p_{i,k}^2}{2m_i}
    +
    V(\vec{q}).
\end{equation}

At thermal equilibrium, the Boltzmann distribution gives the probability of a configuration:
\begin{equation}
    P(\vec{q})
    =
    \frac{1}{Z}
    \exp
    \left(
        -\frac{V(\vec{q})}{k_B T}
    \right),
\end{equation}
with partition function
\begin{equation}
    Z
    =
    \int
    \exp
    \left(
        -\frac{V(\vec{q})}{k_B T}
    \right)
    d\vec{q}.
\end{equation}
Taking a gradient with respect to $q_{i,k}$ gives
\begin{align}
    \frac{dp_{i,k}}{dt}
    &=
    -
    \frac{\partial H}{\partial q_{i,k}}
    =
    -
    \frac{\partial}{\partial q_{i,k}}
    V(\vec{q})
    \nonumber \\
    &=
    k_B T
    \frac{\partial}{\partial q_{i,k}}
    \log P(\vec{q}).
\end{align}
Thus, if a model can evaluate or differentiate the equilibrium log density, it can provide a force-like quantity.
This relation bridges generative density modeling and trajectory inference.
In the present codebase, we test this idea first in the overdamped Brownian setting, where the corresponding relation is $F(x)=D\nabla_x\log p_{\mathrm{eq}}(x)$, or equivalently $\nabla_x\log p_{\mathrm{eq}}(x)=F(x)/D$.

\section{Data}

The training data contain centered bead configurations from simulated Rouse trajectories.
The main data path loads HDF5 files containing arrays of shape
\begin{equation}
    x \in \mathbb{R}^{d \times N \times B},
\end{equation}
where $B$ denotes the number of frames.
The current polymer configs use two-dimensional chains with $N=32$ beads.
For the base Rouse simulation, the SDE above generates trajectories with independent Brownian noise in every bead coordinate.
The base simulation uses $D=0.25$, unit bond length, and $k_\xi=3.0$.
The raw trajectories may include center-of-mass diffusion, but model training centers each frame so the models learn internal conformations rather than global translation.

\section{Time Conventions}

The manuscript uses three time variables with different meanings.
In the Brownian dynamics simulator, $t$ represents physical simulation time.
It indexes the stochastic trajectory that generates equilibrium snapshots.
In the CNF, we use $\tau \in [0,1]$ as a model pseudotime: $\tau=0$ is the data distribution and $\tau=1$ is the standard normal base distribution.
Likelihood evaluation integrates from $\tau=0$ to $\tau=1$, while sampling integrates back from $\tau=1$ to $\tau=0$.
In the diffusion model, $t \in [0,1]$ represents noising time: $t=0$ gives clean data and $t=1$ gives approximately standard Gaussian noise.
The reverse diffusion sampler starts at $t=1$ and integrates to a small positive time, currently $10^{-4}$, for numerical stability.
These conventions match the implementation.

\section{Synthetic Systems}

We focus on two controlled synthetic systems: the base Rouse chain and a hairpin-stabilized Rouse chain.
The base system gives the simplest test of whether a model can learn the equilibrium distribution and score of a flexible Brownian polymer.
The hairpin system adds local structure through attractive nonlocal contacts and repulsive excluded volume.
We do not intend it as a protein model, but it gives a minimal analogue of a stable folded motif: a chain has flexible thermal motion, but a subset of long-range contacts stabilizes a preferred conformation.

\subsection*{Base Rouse Chain}

The base Rouse model uses a quadratic bond potential between adjacent beads.
In the notation of the score model, we write the dimensionless equilibrium potential as
\begin{equation}
    U_{\mathrm{bond}}(x)
    =
    \frac{k_\xi}{2D}
    \sum_{i=2}^{N}
    \|x_i - x_{i-1}\|^2.
\end{equation}
The corresponding equilibrium score satisfies
\begin{equation}
    \nabla_x \log p_{\mathrm{eq}}(x)
    =
    -
    \nabla_x U_{\mathrm{bond}}(x),
\end{equation}
and the Brownian drift satisfies $F(x)=D\nabla_x\log p_{\mathrm{eq}}(x)$.
The base configuration uses $N=32$ beads in two dimensions, $D=0.25$, and $k_\xi=3.0$.
This setting disables all nonideal interactions.
The base Rouse chain therefore provides a clean baseline because adjacent harmonic springs generate the equilibrium score entirely.

\subsection*{Nonideal Interactions}

The simulator also supports several nonideal potentials.
For nonbonded Lennard-Jones interactions, the potential has the form
\begin{equation}
    U_{\mathrm{LJ}}(r)
    =
    4\epsilon
    \left[
        \left(\frac{\sigma}{r}\right)^{12}
        -
        \left(\frac{\sigma}{r}\right)^6
    \right],
\end{equation}
with optional cutoff, softening, bonded-neighbor exclusion, and energy shifting at the cutoff.
The force implementation uses the corresponding radial derivative with softened squared distance $r^2+\delta^2$ for numerical stability.

Excluded volume uses a purely repulsive power-law potential,
\begin{equation}
    U_{\mathrm{EV}}(r)
    =
    \epsilon_{\mathrm{EV}}
    \left(
        \frac{\sigma_{\mathrm{EV}}}{r}
    \right)^{p},
\end{equation}
again with cutoff, softening, and bonded-neighbor exclusion.
The current excluded-volume configuration inherits these defaults from the base file:
\begin{align}
    \epsilon_{\mathrm{EV}}=0.03,
    \quad
    \sigma_{\mathrm{EV}}=1.0,
    \quad
    p=12,
    \nonumber \\
    \delta=0.3,
    \quad
    r_{\mathrm{cut}}=1.5.
\end{align}
This term prevents physically implausible bead overlap while preserving the stochastic flexibility of the chain.

\subsection*{Hairpin-Stabilized Chain}

The hairpin system starts from a two-strand initialization.
For an even number of beads, the initializer places the first half of the chain on one strand and the second half on the opposing strand, with paired beads facing each other.
The current hairpin configuration uses an initial strand separation of $0.95$ and Gaussian initialization noise with standard deviation $0.02$.

Attractive Lennard-Jones contacts between paired beads stabilize the hairpin:
\begin{equation}
    i+j=N+1,
    \qquad
    |i-j| \geq s_{\min}.
\end{equation}
Only these designated long-range pairs receive the hairpin LJ interaction.
The current configuration uses
\begin{align}
    \epsilon_{\mathrm{hp}}=5.0,
    \quad
    \sigma_{\mathrm{hp}}=0.9,
    \quad
    \delta_{\mathrm{hp}}=0.3,
    \nonumber \\
    r_{\mathrm{cut}}=3.0,
    \quad
    s_{\min}=2.
\end{align}
The hairpin configuration also turns on excluded volume and centered confinement.
The confinement term uses the potential
\begin{equation}
    U_{\mathrm{conf}}(x)
    =
    c
    \sum_{i=1}^{N}
    \|x_i-\bar{x}\|^2,
    \qquad
    c=0.02,
\end{equation}
where $\bar{x}$ denotes the chain center of mass.
Together, these terms create a stable but thermally fluctuating structure.
This synthetic example comes closest to the protein motivation: the model must learn not only a broad polymer ensemble, but also a distribution shaped by stabilizing nonlocal contacts.

\section{Continuous Normalizing Flow}

The CNF model learns an ODE that maps data coordinates to a standard normal base distribution.
Let $z(\tau)$ denote the CNF state at pseudotime $\tau \in [0,1]$:
\begin{equation}
    \frac{dz(\tau)}{d\tau} = f_\theta(z(\tau), \tau).
\end{equation}
Here, $\tau$ is a generative-model coordinate, not physical Brownian time.
The implementation integrates the CNF from data to noise for likelihood evaluation and from noise to data for sampling.

The log-density evolves according to the instantaneous change-of-variables formula:
\begin{equation}
    \frac{d}{d\tau}\log p(z(\tau))
    =
    -\operatorname{Tr}
    \left(
        \frac{\partial f_\theta}{\partial z}
    \right).
\end{equation}
Equivalently, the accumulated log-density change satisfies
\begin{align}
    \log P(z(\tau)) - \log P(z(0))
    &=
    \Delta \log P_z(\tau)
    \nonumber \\
    &=
    \int_0^\tau
    -
    \operatorname{Tr}
    \left(
        \frac{\partial f_\theta}{\partial z(\lambda)}
    \right)
    d\lambda.
\end{align}
For a training sample $x$, the model solves
\begin{align}
    z(0) &= x, \\
    z(1) &= \mathrm{ODESolve}(z(0), f_\theta, 0, 1).
\end{align}
The model then computes the log likelihood as
\begin{equation}
    \log p_\theta(x)
    =
    \log \mathcal{N}(z(1); 0, I)
    -
    \Delta \log p,
    \end{equation}
where the ODE solve accumulates $\Delta \log p$.
We can write the same calculation as the augmented ODE system
\begin{align}
    \begin{bmatrix}
        z(1) \\
        \Delta \log P_z(1)
    \end{bmatrix}
    &=
    \begin{bmatrix}
        z(0)=x \\
        \Delta \log P_z(0)=0
    \end{bmatrix}
    \nonumber \\
    &\quad
    +
    \int_0^1
    \begin{bmatrix}
        f_\theta(z(\tau),\tau) \\
        -\operatorname{Tr}
        \left(
            \frac{\partial f_\theta}{\partial z(\tau)}
        \right)
    \end{bmatrix}
    d\tau.
\end{align}
The training objective minimizes negative log likelihood:
\begin{equation}
    \mathcal{L}_{\mathrm{CNF}}(\theta)
    =
    -\mathbb{E}_{x \sim p_{\mathrm{data}}}
    \left[
        \log p_\theta(x)
    \right].
\end{equation}

The current polymer CNF uses an EGNN-style vector field.
The network builds pairwise messages from node embeddings, sequence separation $i-j$, squared pair distances $\|x_i-x_j\|^2$, and pseudotime $\tau$.
The network forms coordinate updates from relative displacements, so the vector field remains translation equivariant after centering.
The main EGNN CNF configuration estimates the divergence term with a Hutchinson JVP trace estimator rather than an exact trace.
Sampling starts from Gaussian noise and integrates the same CNF dynamics backward from $\tau=1$ to $\tau=0$.

\subsection*{Adjoint Sensitivity Equations}

The original CNF formulation can compute parameter gradients with an adjoint ODE.
For the scalar objective
\begin{equation}
    \mathcal{L}
    =
    \log P(z(1)) - \Delta \log P_z(1),
\end{equation}
define the adjoint state
\begin{equation}
    a(\tau)
    =
    \begin{bmatrix}
        a_z(\tau) \\
        a_\Delta(\tau)
    \end{bmatrix}
    =
    \begin{bmatrix}
        \partial \mathcal{L}/\partial z \\
        \partial \mathcal{L}/\partial(\Delta \log P_z)
    \end{bmatrix}.
\end{equation}
At $\tau=1$,
\begin{equation}
    a(1)
    =
    \begin{bmatrix}
        \nabla_z \log P(z(1)) \\
        -1
    \end{bmatrix}.
\end{equation}
The adjoint dynamics satisfy
\begin{align}
    \frac{da(\tau)}{d\tau}
    &=
    -
    \begin{bmatrix}
        \frac{\partial f_\theta}{\partial z} & 0 \\
        -\nabla_z
        \operatorname{Tr}
        \left(
            \frac{\partial f_\theta}{\partial z}
        \right)
        & 0
    \end{bmatrix}^{T}
    a(\tau),
    \\
    \frac{d}{d\tau}
    \left(
        \frac{\partial \mathcal{L}}{\partial \theta}
    \right)
    &=
    a(\tau)^T
    \begin{bmatrix}
        \frac{\partial f_\theta}{\partial \theta} \\
        -\nabla_\theta
        \operatorname{Tr}
        \left(
            \frac{\partial f_\theta}{\partial z}
        \right)
    \end{bmatrix}.
\end{align}
The Julia implementation uses differentiable ODE solves and adjoint machinery through the SciML stack rather than manually coding these equations in the experiment loop.

\section{Variance-Preserving Diffusion Model}

The diffusion baseline learns the score field directly.
The implementation uses a continuous-time variance-preserving (VP) forward SDE
\begin{equation}
    dx_t
    =
    -\frac{1}{2}\beta(t)x_t\,dt
    +
    \sqrt{\beta(t)}\,dw_t,
\end{equation}
with a linear noise-rate schedule
\begin{equation}
    \beta(t) = \beta_{\min} + t(\beta_{\max}-\beta_{\min}),
    \quad
    t \in [0,1],
\end{equation}
with $\beta_{\min}=0.1$ and $\beta_{\max}=20$ in the code.
The marginal forward process has the form
\begin{equation}
    x_t = \mu(t)x_0 + \sigma(t)\epsilon,
    \quad
    \epsilon \sim \mathcal{N}(0,I),
\end{equation}
where
\begin{align}
    \mu(t)
    &=
    \exp
    \left[
        -\frac{1}{4}(\beta_{\max}-\beta_{\min})t^2
        -
        \frac{1}{2}\beta_{\min}t
    \right], \\
    \sigma^2(t)
    &=
    1
    -
    \exp
    \left[
        -\frac{1}{2}(\beta_{\max}-\beta_{\min})t^2
        -
        \beta_{\min}t
    \right].
\end{align}
The implementation centers both $x_0$ and the sampled noise, so the diffusion process stays in the centered coordinate subspace that training uses.
For numerical stability, the implementation clamps training times below by $10^{-5}$ and floors $\sigma(t)$ at the corresponding minimum positive value before forming the score target.

We train the score network $s_\theta(x_t,t)$ to estimate
\begin{equation}
    \nabla_{x_t}\log p_{0t}(x_t|x_0)
    =
    -\frac{\epsilon}{\sigma(t)}.
\end{equation}
The implementation uses the weighted denoising score-matching objective
\begin{equation}
    \mathcal{L}_{\mathrm{diff}}(\theta)
    =
    \mathbb{E}_{t,x_0,\epsilon}
    \left[
        \sigma^2(t)
        \left\|
            s_\theta(x_t,t)
            +
            \frac{\epsilon}{\sigma(t)}
        \right\|_2^2
    \right].
\end{equation}

The diffusion score model also uses an EGNN-style architecture.
It keeps learned node embeddings, uses sequence-distance and radial pair features, and updates coordinates through pairwise relative displacements.
The implementation centers the output and interprets it directly as $\nabla_x \log p_t(x)$.

Generation uses the reverse-time VP SDE.
For this linear schedule, the reverse coefficients satisfy
\begin{equation}
    f(t) = -\frac{1}{2}\beta(t),
    \quad
    g(t)^2 = \beta(t).
\end{equation}
The Euler-Maruyama sampler starts from centered Gaussian noise at $t=1$ and steps backward to a small positive time:
\begin{align}
    x_{t-\Delta t}
    &=
    x_t
    -
    \left[
        f(t)x_t
        -
        g(t)^2 s_\theta(x_t,t)
    \right]\Delta t
    \nonumber \\
    &\quad
    +
    g(t)\sqrt{\Delta t}\,z,
    \qquad
    z \sim \mathcal{N}(0,I).
\end{align}
The sampler integrates from $t=1$ to $t=10^{-4}$ by default.
For numerical stability, it clips sample score norms to $64$, recenters each state, checks for non-finite samples, and omits the final noise draw on the last step.
During diffusion training, the optimizer also clips parameter gradients elementwise to magnitude $1$.

\section{Evaluation}

The experiments report sample-level statistics and score-field diagnostics.
For generated samples, the experiments use the mean absolute error between average pairwise distances in generated configurations and training configurations as the main structural metric.
The code also records the fraction of generated values and samples that remain finite.

For Rouse and polymer Langevin data, the simulator provides analytic force targets.
We obtain the CNF score by differentiating the learned log likelihood,
\begin{equation}
    \nabla_x \log p_\theta(x),
\end{equation}
while we take the diffusion score from the network output at a small time value near zero.
The evaluation scripts compare these scores with the known Brownian score target $F(x)/D$ and apply normalization corrections when the data normalizer runs.
This makes the benchmark more direct than visual sample comparison alone: the model should not only generate plausible chain conformations, but also recover the local equilibrium score implied by the Brownian dynamics.

\section{Julia and SciML Implementation}

We build the implementation around the SciML ecosystem.
This choice does more than provide convenience: the project needs stochastic simulation, neural ODE likelihoods, differentiable ODE solves, GPU-compatible array code, and several trace-estimation strategies in the same code path.
Julia makes these pieces composable without translating between separate simulation, autodiff, and machine-learning frameworks.

\subsection*{Brownian Data Generation with StochasticDiffEq}

The Rouse data generator constructs the simulation directly as an SDE problem.
The simulator constructs
\begin{equation}
    dX_t = F(X_t)\,dt + \sqrt{2D}\,dW_t
\end{equation}
as a \texttt{SDEProblem}; the drift function calls the same polymer force routine that the evaluation code later uses for score targets.
The raw trajectory script uses \texttt{StochasticDiffEq} solvers through a configurable solver interface.
The available choices include Euler-Maruyama, Euler-Heun, RKMil, SOSRA, and SOSRA2.
This lets the same model family generate inexpensive baseline trajectories or more stable stochastic trajectories by changing a YAML configuration rather than rewriting simulation code.

The learning experiments then reuse the saved trajectories through HDF5 loaders.
This structure separates stochastic simulation from generative-model training while keeping both stages inside Julia.

\subsection*{CNF Likelihoods as SciML ODE Problems}

The CNF implementation defines an \texttt{ODEProblem} whose \texttt{ComponentArray} state contains both coordinates and accumulated log-density change:
\begin{equation}
    u(\tau)
    =
    \left(
        x(\tau),
        \Delta \log p(\tau)
    \right).
\end{equation}
The ODE right-hand side returns
\begin{equation}
    \frac{d}{d\tau}
    \begin{bmatrix}
        x(\tau) \\
        \Delta \log p(\tau)
    \end{bmatrix}
    =
    \begin{bmatrix}
        f_\theta(x(\tau),\tau) \\
        -\operatorname{Tr}
        \left(
            \partial f_\theta/\partial x
        \right)
    \end{bmatrix}.
\end{equation}
The code solves the same problem definition forward for likelihood evaluation and backward for sample generation.
The configs select the solver from the SciML ODE stack; current configs use adaptive \texttt{Tsit5} and support fixed-step \texttt{RK4}.
For training, the solve uses \texttt{InterpolatingAdjoint(autojacvec=ZygoteVJP())}, so gradients propagate through the ODE solve while remaining compatible with the neural vector field and the trace estimator.

This design keeps the generative model close to its mathematical definition.
The code does not need a separate CNF framework: the likelihood model is an ODE with a structured state, a vector field, a trace estimator, and a SciML adjoint.

\subsection*{Hutchinson JVP Trace Estimator}

The main EGNN CNF configuration uses a Hutchinson trace estimator with a Jacobian-vector product (JVP).
For a vector field $f_\theta(x,\tau)$, the Hutchinson identity estimates the divergence as
\begin{equation}
    \operatorname{Tr}
    \left(
        \frac{\partial f_\theta}{\partial x}
    \right)
    =
    \mathbb{E}_{v}
    \left[
        v^T
        \frac{\partial f_\theta}{\partial x}
        v
    \right],
\end{equation}
where $v$ denotes a Rademacher probe.
The implementation samples one probe per batch and computes
\begin{equation}
    \widehat{\operatorname{Tr}}
    =
    \sum_{i,j}
    v_{i,j}
    \left[
        J_{f_\theta}(x,\tau)v
    \right]_{i,j}.
\end{equation}

The implementation propagates the EGNN JVP explicitly through the vector field.
The implementation carries both the primal coordinate update and the directional derivative through pairwise displacements, squared distances, edge networks, node updates, and coordinate aggregation.
For example, if $r_{ij}=x_i-x_j$, then the JVP carries the directional derivative
\begin{equation}
    dr_{ij}=v_i-v_j,
    \qquad
    d\|r_{ij}\|^2
    =
    2r_{ij}^{T}dr_{ij}.
\end{equation}
The implementation passes these directional features through the same edge and node networks as the primal features, producing both $f_\theta(x,\tau)$ and $J_{f_\theta}(x,\tau)v$ in one structured forward pass.
This avoids materializing the full Jacobian over all bead coordinates.
For a 32-bead chain in two dimensions, the full Jacobian already forms a dense object over 64 coordinate variables per sample; for larger systems, exact traces become the wrong abstraction.
The JVP estimator keeps the CNF likelihood scalable while preserving the same ODE interface.

The code also keeps alternative trace modes and provides an exact Zygote-based divergence for smaller checks.
It also provides a finite-difference Hutchinson mode, \texttt{hutchinson\_fd}, which approximates the same JVP by
\begin{equation}
    J_{f_\theta}(x,\tau)v
    \approx
    \frac{
        f_\theta(x+\epsilon v,\tau)
        -
        f_\theta(x-\epsilon v,\tau)
    }{2\epsilon}.
\end{equation}
This finite-difference path provides a useful diagnostic fallback because it tests the same trace estimator without relying on the analytic JVP implementation.
The production path uses \texttt{hutchinson\_jvp} because it runs faster and remains analytic: it avoids two extra vector-field evaluations per trace estimate, avoids finite-difference step-size error, and remains differentiable through the exact directional-derivative computation.
Training still uses Zygote for gradients through the CNF loss and SciML adjoint solve, but Zygote does not derive this EGNN JVP automatically; the model code implements the JVP explicitly.

\subsection*{Autodiff and Stability}

Both CNF and diffusion training use Zygote for reverse-mode gradients.
Several small custom adjoints remove avoidable AD overhead for common tensor operations such as batched node embeddings, row concatenation, and centering.
The diffusion implementation also uses \texttt{Zygote.ignore} and \texttt{dropgrad} where fixed geometric inputs, masks, sampled noise, or numerical guards should not become part of the trainable derivative path.

The result keeps stochastic simulation, adaptive ODE solving, adjoint differentiation, neural network training, and GPU execution in one language and one type system.
This gives Julia its main advantage for this project: the mathematical objects in the paper map directly to executable SciML problems, and the engineering choices needed for stability remain visible in the code rather than hidden behind a separate modeling layer.

\section{Trajectory-Inference Algorithm}

We can summarize the full long-term BoltzFlow trajectory-inference idea from the earlier proposal as follows.
First, obtain equilibrium snapshots
\begin{equation}
    \{\vec{q}_1,\ldots,\vec{q}_n\}
    \sim
    P(\vec{q}).
\end{equation}
Then train a CNF to model $P(\vec{q})$ and evaluate $\log P(\vec{q})$.
For a trajectory initialized at $\vec{q}(t_0)$, sample momenta from the Maxwell-Boltzmann distribution,
\begin{equation}
    p_{i,k}(t_0)
    \sim
    \mathcal{N}
    \left(
        0,
        m_i k_B T
    \right).
\end{equation}
At each physical timestep, compute the CNF log density and its gradient:
\begin{equation}
    \nabla_{\vec{q}} \log P(\vec{q}(t)).
\end{equation}
Using the Boltzmann relation,
\begin{equation}
    -
    \nabla_{\vec{q}} V(\vec{q})
    =
    k_B T
    \nabla_{\vec{q}}
    \log P(\vec{q}),
\end{equation}
update the momenta and positions with Hamilton's equations:
\begin{align}
    \frac{d\vec{p}}{dt}
    &=
    k_B T
    \nabla_{\vec{q}}
    \log P(\vec{q}),
    \\
    \frac{d\vec{q}}{dt}
    &=
    \frac{\partial H}{\partial \vec{p}}.
\end{align}
This Hamiltonian trajectory-inference loop provides the protein-facing motivation.
The current implementation does not yet execute this loop for atomistic proteins; it instead validates the score-learning components on Brownian/Rouse systems where we know the target force field.

\section{Current Scope}

We should describe the current codebase as a Brownian-dynamics and polymer-equilibrium study.
Protein dynamics remain a longer-term motivation, but the implemented system does not yet infer atomistic protein trajectories from static structural ensembles.
The code supports a narrower practical claim: we can train equivariant CNF and diffusion models on centered equilibrium snapshots from Rouse-style Brownian simulations, sample them as generative models, and evaluate them against known equilibrium score fields.

% \bibliographystyle{apsrev4-2}
% \bibliography{references}

\end{document}
